% ============================================================================
% Fichier  : front/02_mode_emploi.tex
% Titre    : Mode d'emploi du cours
% Auteur   : MINKA MI NGUIDJOÏ Thierry Emmanuel
% Version  : 1.0 -- novembre 2026
% ============================================================================
% 60 h = 30 séances de 2 h (norme des cours ENSPY) ; fiches et labs hors
% de ces 60 h.
% ============================================================================

\chapter*{Mode d'emploi du cours}
\addcontentsline{toc}{chapter}{Mode d'emploi du cours}
\markboth{Mode d'emploi}{Mode d'emploi}

\section*{Structure générale}

Le cours compte un prologue et six parties, soit 31~chapitres, pour
60~heures en présentiel réparties en trente séances de deux heures. Les
fiches d'auto-évaluation par chapitre et les laboratoires de 60~minutes de
NDAP SUKULU sont des travaux à faire hors de ces 60~heures. Les Parties~V
et~VI servent en outre de point de départ aux projets et aux mémoires.

\begin{tableaucours}{L{1.6cm}L{4.4cm}L{4.4cm}L{1.8cm}}{%
    \tableauHeader{Partie} &
    \tableauHeader{Thème} &
    \tableauHeader{Compétences visées} &
    \tableauHeader{Niveau}
}
Prologue & Quarante-cinq questions & Distinguer constatation, impression et
opinion & \niveaubase \\
I & Fondements de l'audit, du contrôle et du risque & Situer l'audit,
évaluer un contrôle, mesurer et gérer un risque, repérer la fraude &
\niveaubase \\
II & Méthode de la mission & Planifier, décomposer, collecter la preuve,
conclure avec $\mu$, $U$ et $H$, suivre & \niveaumoyen \\
III & Cadre normatif, juridique et déontologique & Choisir ses normes et
ses critères~; connaître ses obligations et ses devoirs & \niveaubase \\
IV & Pratique de l'audit des SI & Auditer la gouvernance, une application,
un projet, la sécurité, la continuité, les tiers, la conformité &
\niveaumoyen{} à \niveauexpert \\
V & Défis avancés & Audit continu, IA, passage à la preuve judiciaire &
\niveauexpert \\
VI & Recherche~: l'auditologie & Lire et prolonger le programme de
recherche & \niveauexpert \\
\end{tableaucours}

\section*{Les trente séances}

\begin{tableaucours}{L{1.4cm}L{7.4cm}L{3.4cm}}{%
    \tableauHeader{Séance} &
    \tableauHeader{Contenu} &
    \tableauHeader{Partie}
}
1 & Prologue~: quarante-cinq questions & Prologue \\
2 & Auditer, contrôler, donner une assurance (Ch.~\ref{chap:audit-auditologie}) & I \\
3 & Grandes défaillances (Ch.~\ref{chap:defaillances}) & I \\
4 & Le contrôle interne (Ch.~\ref{chap:controle-interne}) & I \\
5 & Le risque, notions et mesure (Ch.~\ref{chap:risque}) & I \\
6 & Gérer les risques des SI (Ch.~\ref{chap:gestion-risques}) & I \\
7 & La fraude (Ch.~\ref{chap:fraude}) & I \\
8 & La mission d'audit (Ch.~\ref{chap:mission}) & II \\
9 & Préparer et planifier (Ch.~\ref{chap:planifier}) & II \\
10 & Décomposer l'espace d'audit (Ch.~\ref{chap:decomposition}) & II \\
11 & Collecter la preuve (Ch.~\ref{chap:preuve}) & II \\
12 & Conclure et rédiger (Ch.~\ref{chap:conclure}) & II \\
13 & Suivre (Ch.~\ref{chap:suivi}) & II \\
14 & Les normes professionnelles (Ch.~\ref{chap:normes}) & III \\
15 & Les référentiels de contrôle (Ch.~\ref{chap:referentiels}) & III \\
16 & Le cadre juridique (Ch.~\ref{chap:droit}) & III \\
17 & Déontologie et profession (Ch.~\ref{chap:deontologie}) & III \\
18 & Gouvernance et fonction informatique (Ch.~\ref{chap:audit-gouvernance}) & IV \\
19 & Une application en service (Ch.~\ref{chap:audit-applications}) & IV \\
20 & Les projets informatiques (Ch.~\ref{chap:audit-projets}) & IV \\
21 & Le système de management de la sécurité (Ch.~\ref{chap:audit-smsi}) & IV \\
22 & La sécurité technique (Ch.~\ref{chap:audit-technique}) & IV \\
23 & Tests d'intrusion et audit de code (Ch.~\ref{chap:intrusion}) & IV \\
24 & Production et continuité (Ch.~\ref{chap:audit-continuite}) & IV \\
25 & Les tiers et le nuage (Ch.~\ref{chap:audit-tiers}) & IV \\
26 & La conformité (Ch.~\ref{chap:audit-conformite}) & IV \\
27 & Études de cas intégratrices (Ch.~\ref{chap:cas-integrateurs}) & IV \\
28 & Audit continu~; auditer l'IA (Ch.~\ref{chap:audit-continu}
et~\ref{chap:audit-ia}) & V \\
29 & De l'audit à la preuve judiciaire (Ch.~\ref{chap:audit-preuve}) & V \\
30 & L'auditologie~; mesurer l'incertitude (Ch.~\ref{chap:auditologie}
et~\ref{chap:incertitude}) & VI \\
\end{tableaucours}

\section*{Les encadrés}

\begin{boxdef}
\textbf{Définition}. Introduit un terme avec sa définition~; les termes
définis sont indexés et repris dans le glossaire (Annexe~\ref{ann:glossaire}).
\end{boxdef}

\begin{boxscenario}{Titre du scénario}
\textbf{Scénario opérationnel}. Un épisode d'une des trois opérations du
cours~: LIZBIZ, SANAGA, MEFOU.
\end{boxscenario}

\begin{boxlecture}{Titre de la lecture}
\textbf{Lecture auditologique}. Une cotation \OUI{}~/ \NON{}~/ \NEK{}, sa
constatation, son analyse avec un niveau de confiance, et les trois mesures
du domaine, rappelées en bas à droite.

\medskip
\hfill\lecture{0{,}65}{0{,}25}{1{,}24}
\end{boxlecture}

\begin{boxobjet}{Code · Intitulé}
\textbf{Fiche d'assertion}. Présente un contrôle sous une forme
constante~: objectif, critère, question, risque, conséquence, puis test,
preuve et trace~; ou une observation en sept éléments.
\end{boxobjet}

\begin{boxprocedure}
\textbf{Papier de travail}. Dit ce que l'auditeur doit faire, et ce que le
dossier doit contenir, pour qu'une étape soit opposable.
\end{boxprocedure}

\begin{boxmetier}{Auditeur SI}
\textbf{Application métier}. Traduit le chapitre en gestes concrets pour
l'un des métiers visés~: auditeur SI, contrôleur interne, auditeur de
sécurité, RSSI, gestionnaire de risques, responsable de la conformité.
\end{boxmetier}

\begin{boxoutils}{Titre de l'outil}
\textbf{Outils}. Un programme, une requête ou une commande, à refaire sur
les données du cas.
\end{boxoutils}

\begin{boxjuris}
Cite un texte de loi ou de règlement, avec son degré de vérification, puis
en tire la portée pour l'audit.
\end{boxjuris}

\begin{boxalert}
\textbf{Violation typique} et \textbf{règle}~: une erreur fréquente, et ce
qu'il faut faire à la place.
\end{boxalert}

\begin{boxresume}
\textbf{À retenir}. Les points essentiels du chapitre.
\end{boxresume}

\begin{boxexo}
\textbf{Exercice}. Dix par chapitre, gradués en trois niveaux~:
\niveaubase\ fondamental, \niveaumoyen\ intermédiaire, \niveauexpert\
expert. Les corrigés figurent dans le livret de correction.
\end{boxexo}

\section*{Conventions}

\begin{itemize}
    \item Les réponses s'écrivent \OUI{}, \NON{} et \NEK{}~; un contrôle sans
          objet est écarté, raison écrite, et n'entre pas dans les calculs.
    \item $\mu$ est le taux de conformité, $U$ le taux d'incertitude, $H$
          l'entropie en bits (au plus $\log_2 3 \approx 1{,}585$)~; la règle
          de décision examine $U$ d'abord, $\mu$ ensuite, $H$ pour contrôler.
    \item Les seuils du quadrant ($\mu = 0{,}5$~; $U = 0{,}20$) et la
          conversion en niveaux de maturité sont des conventions
          pédagogiques, déclarées comme telles.
    \item Les axiomes de l'auditologie portent la numérotation de l'article
          A0 (Chapitre~\ref{chap:auditologie}).
    \item Chaque analyse indique un niveau de confiance~: faible, moyenne ou
          élevée.
    \item Les \termecle{termes clés} apparaissent en gras~; les normes
          s'écrivent \norme{ISO/IEC~27001}, les lois camerounaises
          \loicam{n°~2010/012}.
\end{itemize}

\section*{Matériel d'accompagnement}

\begin{itemize}
    \item \textbf{Annexes}~: glossaire, référentiel des outils, modèles et
          grilles, textes juridiques et normatifs, bibliographie commentée,
          contacts (Annexes~\ref{ann:glossaire} à~\ref{ann:contacts}).
    \item \textbf{Livret de correction}~: corrigés des exercices et des
          fiches, réservé au formateur.
    \item \textbf{NDAP SUKULU}~: fiches d'auto-évaluation et laboratoires,
          présentés dans la section qui suit ce mode d'emploi.
\end{itemize}
